% Irodov Part V §5.7, originals 281--292 (end of Part V).
% Supplied Arihant PDF pp.240 / printed p.231; Appendix and answers pp.347-348.
% In problem 292, photon gains energy because the incident electron STOPS:
% the signed wavelength displacement is negative; book gives its magnitude.
\thermalopticsproblem{281}{231}
应用能量、动量守恒证明：一个在真空中自由运动的
电子不可能完全吸收一个真实光子。
\exsolution
在电子吸收光子前的静止参考系，
电子四动量为 $P^\mu=(m_ec,\boldsymbol0)$，
光子四动量为 $K^\mu=(E/c,\boldsymbol k)$，
其中 $E=pc>0$、$K^\mu K_\mu=0$。
若完全吸收，末态仅一个电子，其四动量为
$P_f^\mu=P^\mu+K^\mu$，应满足电子壳条件
$P_f^\mu P_{f\mu}=m_e^2c^2$。
但直接平方可得
\[
 P_f^2=(P+K)^2=P^2+2P\!\cdot\! K+K^2
 =m_e^2c^2+2m_eE>m_e^2c^2,
\]
与自由电子质量不变相矛盾。
故
\[
 \boxed{\text{自由电子不能单独完全吸收真实光子}}.
\]
\excheck 若有晶格、原子核或其他粒子共同承担反冲，
则不受此两体末态限制；本证明使用自由电子这一必要条件。

\thermalopticsproblem{282}{231}
解释康普顿散射的三个现象：
(a) 散射波长增加量与材料种类无关；
(b) 散射角越大、原子序数越低，
移位谱线的强度通常越高；
(c) 散射光中为何还保有未发生位移的谱线？
\exsolution
(a) 对与\emph{近似自由}电子的一次康普顿散射，
能量、动量守恒给
\[
 \boxed{\Delta\lambda=
 \frac{h}{m_ec}(1-\cos\theta)}.
\]
其中仅含电子质量、常数以及散射角，
故不直接依赖散射材料。
(b) 在一定能区内，高散射角意味着更大的能量转移，
更容易使结合电子表现出准自由电子行为；
低原子序数材料中电子结合较弱，
因此可观察到的移位分量比例往往较大。
(c) 被原子强束缚的电子或整个原子／晶格参与
弹性（相干）散射，较大的反冲体系质量使
波长变化很小，形成未位移的瑞利型散射分量。
\excheck (b) 涉及实验能区与电子结构，不能由
康普顿单电子公式孤立地推出普适的单调强度定律。

\thermalopticsproblem{283}{231}
一束窄单色 X 射线在材料中由自由电子散射。
观测散射角为
$\theta_1=60^\circ$ 与 $\theta_2=120^\circ$
的两束散射光波长之比为
$\eta=\lambda_2/\lambda_1=2.0$。
求原入射波长。
\exsolution
记电子康普顿波长
$\lambda_C=h/(m_ec)=2.426\,\mathrm{pm}$。
由康普顿公式，
\[
 \lambda_1=\lambda_0+\lambda_C(1-\cos60^\circ)
 =\lambda_0+\frac12\lambda_C,
\]
\[
 \lambda_2=\lambda_0+\lambda_C(1-\cos120^\circ)
 =\lambda_0+\frac32\lambda_C.
\]
由
$(\lambda_0+\tfrac32\lambda_C)/
(\lambda_0+\tfrac12\lambda_C)=2$
得
\[
 \boxed{\lambda_0=\frac{\lambda_C}{2}
 \simeq1.21\,\mathrm{pm}}.
\]
\excheck 大散射角对应较长波长，
故 $\lambda_2>\lambda_1$ 与题设 $\eta=2$ 一致。

\thermalopticsproblem{284}{231}
能量 $E_0=1.00\,\mathrm{MeV}$ 的光子
被静止自由电子康普顿散射，散射后
光子波长增加了入射波长的
$\eta=25\%$。求反冲电子动能。
\exsolution
光子的能量反比于波长，
$\lambda'=(1+\eta)\lambda_0$ 意味着
$E'=E_0/(1+\eta)$。
自由电子初始静止，能量守恒给电子动能
$T_e=E_0-E'$，从而
\[
 \boxed{T_e=E_0\frac{\eta}{1+\eta}
 =1.00\,\mathrm{MeV}\frac{0.25}{1.25}
 =0.20\,\mathrm{MeV}}.
\]
\excheck “波长增大25%”不等于
“光子能量减少25%”，因为二者反比。

\thermalopticsproblem{285}{231}
波长 $\lambda=6.0\,\mathrm{pm}$ 的光子被
一个静止自由电子以 $90^\circ$ 散射。
求：(a) 散射光子角频率；
(b) 电子反冲动能。
\exsolution
当 $\theta=90^\circ$，
康普顿式给 $\lambda'=\lambda+\lambda_C$，
$\lambda_C=2.426\,\mathrm{pm}$。
所以
\[
 \boxed{\omega'=\frac{2\pi c}{\lambda+\lambda_C}
 \simeq2.24\times10^{20}\,\mathrm{s^{-1}}}.
\]
由能量守恒，
\[
 \boxed{T_e=hc\left(\frac1\lambda-
 \frac1{\lambda+\lambda_C}\right)
 \simeq59.5\,\mathrm{keV}\approx60\,\mathrm{keV}}.
\]
\excheck (a) 求的是\emph{角频率}而非 $\nu=c/\lambda'$，
二者相差 $2\pi$。

\thermalopticsproblem{286}{231}
初能量为 $E=250\,\mathrm{keV}$ 的光子，
被静止自由电子以角 $\theta=120^\circ$ 散射。
求散射光子的能量。
\exsolution
康普顿公式
$\lambda'-\lambda=\lambda_C(1-\cos\theta)$，
以 $E=hc/\lambda$ 改写为能量形式
\[
 E'=\frac E{1+
 \dfrac E{m_ec^2}(1-\cos\theta)}.
\]
由于
$1-\cos120^\circ=3/2$、
$m_ec^2=511\,\mathrm{keV}$，
\[
 \boxed{E'=
 \frac{250}{1+(250/511)\times1.5}\,\mathrm{keV}
 \simeq144\,\mathrm{keV}}.
\]
\excheck $E'<E$，且散射角较大时
转移给电子的能量较多。

\thermalopticsproblem{287}{231}
初始光子动量为
$p=1.02\,\mathrm{MeV}/c$，
散射后动量变为
$p'=0.255\,\mathrm{MeV}/c$，
散射体为静止自由电子。
求光子的散射角。
\exsolution
因 $\lambda=h/p$、
$\lambda'=h/p'$，
康普顿位移方程为
\[
 \frac hp'-\frac hp=
 \frac h{m_ec}(1-\cos\theta).
\]
约去 $h$ 得
\[
 \boxed{1-\cos\theta=
 m_ec\left(\frac1{p'}-\frac1p\right)}.
\]
以 $m_ec^2=0.511\,\mathrm{MeV}$ 代入，
\[
 1-\cos\theta=
 0.511\left(\frac1{0.255}-\frac1{1.02}\right)
 \simeq1.503,
\]
数值取有效位数后 $\cos\theta\simeq-0.50$，
故
\[
 \boxed{\theta\simeq120^\circ}.
\]
\excheck 使用 $\mathrm{MeV}/c$ 表达动量时，
$m_ec$ 也须换成相同的动量单位。

\thermalopticsproblem{288}{231}
光子被静止自由电子以
$\theta=120^\circ$ 散射后，
电子动能为 $T=0.45\,\mathrm{MeV}$。
求散射前光子的能量 $E$。
\exsolution
由能量守恒，散射光子能量 $E'=E-T$。
康普顿能量关系
\[
 \frac1{E'}-\frac1E=
 \frac{1-\cos\theta}{m_ec^2}.
\]
于是
\[
 \frac{T}{E(E-T)}
 =\frac{1-\cos\theta}{m_ec^2}.
\]
令 $A=m_ec^2/(1-\cos\theta)$，得
$E^2-TE-AT=0$，取正根：
\[
 \boxed{E=\frac T2
 \left[1+\sqrt{1+
 \frac{4m_ec^2}{T(1-\cos\theta)}}\right]}.
\]
代入 $T=0.45\,\mathrm{MeV}$、
$1-\cos120^\circ=1.5$，
\[
 \boxed{E\simeq0.68\,\mathrm{MeV}}.
\]
\excheck 解的另一根为负，不满足
$E>T>0$ 的物理条件。

\thermalopticsproblem{289}{231}
在自由电子康普顿散射中，
反冲电子的\emph{最大}动能为
$T_{\max}=0.19\,\mathrm{MeV}$。
求入射 X 射线波长。
\exsolution
电子得到最大能量当光子背散射 $\theta=\pi$。
若入射能量为 $E$，散射后能量为
$E'=E/(1+2E/(m_ec^2))$，
所以
\[
 T_{\max}=E-E'=
 \frac{2E^2}{m_ec^2+2E}.
\]
解二次方程得正根
\[
 \boxed{E=
 \frac{T_{\max}+
 \sqrt{T_{\max}^2+2m_ec^2T_{\max}}}{2}}.
\]
随后 $\lambda=hc/E$，数值为
\[
 \boxed{\lambda\simeq3.7\,\mathrm{pm}}.
\]
\excheck 最大反冲对应 $\theta=180^\circ$，
不是 $90^\circ$。

\thermalopticsproblem{290}{231}
能量 $E=0.15\,\mathrm{MeV}$ 的光子
被静止自由电子散射后，波长增加
$\Delta\lambda=3.0\,\mathrm{pm}$。
求反冲\emph{电子}与入射方向夹角 $\phi$。
\exsolution
光子散射角 $\theta$ 先由
$\Delta\lambda=\lambda_C(1-\cos\theta)$ 求出，
$\cos\theta=1-3.0/2.426\simeq-0.237$。
令初、末光子能量为 $E,E'$，
以能量为动量单位 $pc=E$，
$E'=hc/(\lambda+\Delta\lambda)$，
$\lambda=hc/E$，得
$E'\simeq0.110\,\mathrm{MeV}$。
动量守恒的平行、垂直入射方向分量分别为
\[
 p_{e,x}=\frac{E-E'\cos\theta}{c},
 \qquad p_{e,y}=-\frac{E'\sin\theta}{c}.
\]
所以
\[
 \boxed{\tan\phi=
 \frac{E'\sin\theta}{E-E'\cos\theta}},
 \qquad
 \boxed{\phi\simeq31^\circ}.
\]
\excheck 题目求的是\emph{电子}方向，
不等于光子散射角 $\theta\simeq104^\circ$。

\thermalopticsproblem{291}{231}
能量为电子静能的 $\eta=2.0$ 倍的光子，
与静止自由电子正碰并背向散射。
反冲电子随后以速度垂直于均匀磁场
$B=0.12\,\mathrm T$ 运动。
求其轨道曲率半径。
\exsolution
记 $m=m_e$，
$E=\eta mc^2$。
光子完全背散射时
\[
 E'=\frac{E}{1+2\eta},\qquad
 T=E-E'=\frac{2\eta^2}{1+2\eta}mc^2.
\]
电子满足相对论动量--能量关系
$p_ec=\sqrt{T(T+2mc^2)}$。
代入并化简得
\[
 p_e=\frac{2\eta(1+\eta)}{1+2\eta}mc.
\]
垂直磁场中洛伦兹力提供向心加速度，
相对论曲率半径为 $r=p_e/(eB)$，所以
\[
 \boxed{r=\frac{2\eta(1+\eta)m_ec}{(1+2\eta)eB}
 \simeq3.4\,\mathrm{cm}}.
\]
\excheck 应使用\emph{相对论动量}，
不可在 MeV 电子上直接代 $p=mv_{\rm nonrel}$。

\thermalopticsproblem{292}{231}
一个运动的相对论电子与光子碰撞。
光子方向偏转 $\theta=60^\circ$，
而电子在碰撞后\emph{恰好停止}。
求散射光子的康普顿波长变化。
\exsolution
与通常“静止电子受撞击”不同，
此题入射电子在运动而末态电子静止。
令入射、出射光子四动量为 $K,K'$，
末态电子四动量为 $P_f=(m_ec,\boldsymbol0)$。
动量守恒写作
$P_i+K=P_f+K'$；
利用 $P_i^2=P_f^2=m_e^2c^2$、
$K^2=K'^2=0$，
平方后得到
\[
 m_ec^2(E'-E)=EE'(1-\cos\theta).
\]
所以
$\tfrac1E-\tfrac1{E'}=(1-\cos\theta)/(m_ec^2)$。
换成波长，得到带符号的光子波长变化
\[
 \boxed{\lambda'-\lambda=
 -\frac h{m_ec}(1-\cos\theta)}.
\]
在 $\theta=60^\circ$，
\[
 \boxed{\Delta\lambda\simeq-1.21\,\mathrm{pm}},
 \qquad
 \boxed{|\Delta\lambda|\simeq1.2\,\mathrm{pm}}.
\]
\excheck 此题是\emph{反向康普顿能量传递}，
电子停止使光子能量增加、波长缩短；
书末只给波长变化的模，负号须独立标注。
